\thispagestyle{empty}
\chapter{構文解析器生成系の世界}\label{ch05}
\index{こうぶんかいせききせいせいけい@構文解析器生成系}
\index{ぱーさーじぇねれーた@パーサージェネレータ}
\index{ぱーさーこんびねーた@パーサーコンビネータ}
\index{CFG}
\index{PEG}

第4章では現在知られている構文解析手法について、アイデアと概要を説明しました。しかし、構文解析の世界ではよく知られていることなのですが、第4章で説明した手法は毎回プログラマが手で実装する必要はありません。

というのは、CFG（文脈自由文法）やPEG（その類似表記も含む）によって記述された文法定義から特定の構文解析アルゴリズムを用いた構文解析器を生成する\textgt{構文解析器生成系}（Parser
Generator、パーサージェネレータ）というソフトウェアがあるからです。

簡単に言えば、「文法の定義を書いたら、対応したパーサーのコードを自動生成してくれるツール」です。これにより、複雑な構文解析アルゴリズムを知らなくても、文法定義さえ書ければ高性能なパーサーを作れます。

構文解析器生成系でもっとも代表的なものはYacc（ヤック）あるいはその互換実装であるGNU
Bison（バイソン）でしょう。YaccはLALR(1)法を利用したC言語の構文解析器を生成してくれるソフトウェアです。

ちなみに「Yacc」は「Yet Another Compiler
Compiler」（また別のコンパイラコンパイラ）の略で、「コンパイラを作るためのコンパイラ」という意味です。

この章では構文解析器生成系という種類のソフトウェアの背後にあるアイデアからはじまり、PEGから実際に対応する構文解析器を作る方法、多種多様な構文解析器生成系についての紹介を行います。

また、この章の最後ではある意味構文解析器生成系の一種とも言える\textgt{パーサーコンビネータ}の実装方法について踏み込んで説明します。

パーサーコンビネータとは、簡単に言えば「プログラミング言語に備わったファーストクラス関数やオブジェクトを組み合わせてパーサーを構築する方法」です。通常の構文解析器生成系が「文法定義ファイル」から「Javaコード」を生成するのに対し、パーサーコンビネータではJavaのメソッドを組み合わせて直接パーサーを作ります。単なる手書きパーサーと違い、パーサーコンビネータは規則や文法に対応した関数やオブジェクトがライブラリとして備わっているため、再利用性が高く、柔軟な構文解析が可能です。

構文解析器生成系を作るのは骨が折れますが、パーサーコンビネータであれば、いわゆる「ラムダ式」を持つほとんどのプログラミング言語で比較的簡単に実装できます。本書で利用しているJavaでも同様です。というわけで、本章を読めば皆さんもパーサーコンビネータを明日から自前で実装できるようになります。

\section{Dyck言語の文法とPEGによる構文解析器生成}
\index{Dyck言語}
\index{PEG}
\index{こうぶんかいせききせいせい@構文解析器生成}
\index{こーどせいせい@コード生成}
\index{ばっくとらっく@バックトラック}

これまで何度も登場したDyck言語（括弧の対応が取れた文字列を表す言語）は明らかにLL(1)法でもLR(1)法でもPEGによっても解析可能な言語です。しかし、手書きでパーサーを実装しようとすると、退屈な繰り返しコードが多くなりがちです。

実際のところ、Dyck言語を表現する文法があって、構文解析アルゴリズムがPEGということまでわかれば対応するJavaコードを\textgt{機械的に生成する}ことも可能そうに見えます。特に、構文解析はコード量が多いわりには退屈な繰り返しコードが多いものですから、文法からJavaコードを生成できれば劇的に工数を削減できそうです。

このように「文法と構文解析手法が決まれば、後のコードは自動的に決定可能なはずだから、機械に任せてしまおう」という考え方が構文解析器生成系というソフトウェアの背後にあるアイデアです。

早速ですが、以下のようにDyck言語を表す文法（第3章で示した
\texttt{P\ →\ (\ P\ )\ P\ \textbar{}\ ε}
に近いもの）がPEGで与えられたとして、構文解析器を生成する方法を考えてみましょう。

\subsection{Dyck言語のPEG}

PEGの記法では、\texttt{/}は「順序付き選択」を表し、BNFの\texttt{\textbar{}}とは異なり「最初の選択肢がマッチしたらそれを採用し、2番目以降は試さない」という意味になります。

\begin{codescreen}{}
D <- P;
P <- "(" P ")" P / ""; // "" は空文字列ε（イプシロン）を表す
\end{codescreen}

このPEGは\texttt{P} が \texttt{(P)P}
の形であるか、または空文字列であることを示します。PEGでは非終端記号の呼び出しは関数呼び出しとみなすことができますから、この定義に対応するパーサー関数（メソッド）のスケルトンを考えると、次のようなコードのイメージになります。

\begin{codescreen}{}
// 概念的なスケルトン
public boolean parseD() {
    return parseP(); // DはPに委譲
}

public boolean parseP() {
    // P <- "(" P ")" P / "";
    // 最初に "(" P ")" P を試す：
    //   "(" にマッチするか？
    //   マッチしたら、P を再帰的に呼び出す (1回目)
    //   Pの呼び出しが成功したら： 
    //     ")" にマッチするか？
    //     マッチしたら、P を再帰的に呼び出す (2回目)
    //     Pの呼び出しが成功したら：
    //       全体として成功
    // もし途中で失敗したら：
    //   バックトラックして "" を試す
    //   "" は常に成功（空文字列を消費）
    // どちらかが成功すれば parseP は成功
}
\end{codescreen}

次の項ではこのスケルトンを実際のJavaコードに変換して、Dyck言語の構文解析器を生成する方法を見ていきます。

\subsection{Dyck言語の構文解析器を生成する}

PEGからJavaコードへの機械的な変換は、以下の規則に従います。

\begin{itemize}
\item
  \textgt{非終端記号}\texttt{N}： \texttt{parseN()}
  という名前のメソッドを生成
\item
  \textgt{文字列リテラル}\texttt{str}：\texttt{match(str)}
  の呼び出しに変換
\item
  \textgt{連接} \texttt{e1\ e2}：
  \texttt{e1}の中身を再帰的に変換した結果と\texttt{e2}の中身を再帰的に変換した結果を連結
\item
  \textgt{順序付き選択} \texttt{e1\ /\ e2}：
  \texttt{e1}を試すコードと、失敗したらバックトラックするコード、\texttt{e2}を試すコードを連結する

  \begin{itemize}
  \item
    ここではバックトラックのために例外を使う
  \end{itemize}
\item
  \textgt{繰り返し}：ループ構造に変換
\end{itemize}

\begin{codescreen}{}
public class DyckParser {
    private String input;
    private int pos;

    public static class Failure extends RuntimeException {
        public Failure(String message) {
            super(message);
        }
    }
    
    public DyckParser(String input) {
        this.input = input;
        this.pos = 0;
    }
    
    // D <- P
    public void parseD() {
        parseP();
    }
    
    // P <- "(" P ")" P / ""
    public void parseP() {
        // 現在位置を保存（バックトラック用）
        int saved = pos;

        try {
            match("(");  // "(" にマッチするか？
            parseP();  // Pを再帰的に呼び出す（1回目）
            match(")");  // ")" にマッチするか？
            parseP();  // Pを再帰的に呼び出す（2回目）
            return;  // 成功
        } catch (Failure e) {
            // 失敗したらバックトラック
            pos = saved;
            // 空文字列 "" を試す（常に成功）
            return;
        }
    }
    
    // 文字列のマッチングを行うヘルパーメソッド
    private void match(String str) {
        if (pos + str.length() <= input.length() &&
            input.startsWith(str, pos)) {
            pos += str.length();
            return;
        }
        throw new Failure("Expected '" + str + "' at position " + pos);
    }
    
    // パース実行メソッド
    public boolean parse() {
        parseD();
        // 例外が起きず、全ての文字が消費されていれば成功
        if (pos == input.length()) {
            return true;  
        } else {
            return false;
        }
    }
}
\end{codescreen}

このようにして、比較的シンプルにPEGからDyck言語の構文解析器を生成することができました。例外を使ってバックトラックを実装していますが、これは単純化のためです。Javaで例外を投げるのはそこそこ重い処理なので、単純にif文で分岐する方法を使った方が効率的です。

\subsection{構文解析器生成系の実装}

実際の構文解析器生成系では、以下のようなステップで自動生成を行います。

\begin{enumerate}
\item
  \textbf{\sffamily PEG}\textgt{文法の解析}：PEG記法で書かれた文法定義自体を解析
\item
  \textgt{中間表現の生成}：文法規則を内部的なデータ構造に変換
\item
  \textgt{コード生成}：テンプレートを使ってJavaコードを出力
\end{enumerate}

以下は中間表現が生成された後のコード生成の例です。ここでは、PEGの規則を表すクラス\texttt{Rule}と、各種PEG式を表す\texttt{Expr}クラスを定義し、Javaコードを生成しています。

\begin{codescreen}{}
import java.util.*;
import java.util.stream.Collectors;
// 簡単な構文解析器生成系の例
public class PEG2Java {
    // 文法規則を表すレコード
    public record Rule(String name, Expr body) {}
    
    // PEG式を表す抽象クラス
    public abstract static class Expr {
        abstract String generate(int indentLevel);
    }
    
    // 文字列リテラル
    public static class Lit extends Expr {
        public final String value;
        public Lit(String value) {
            this.value = value;
        }
        String generate(int indentLevel) {
            return ("match(\"" + value + "\");").indent(indentLevel);
        }
    }

    // 非終端記号
    public static class NT extends Expr {
        public final String name;
        public NT(String name) {
            this.name = name;
        }
        // 非終端記号のコード生成
        String generate(int indentLevel) {
            return ("parse" + name + "();").indent(indentLevel);
        }
    }
    
    // 連接
    static class Seq extends Expr {
        public final List<Expr> exprs;
        public Seq(Expr... exprs) {
            this.exprs = List.of(exprs);
        }

        // 連接のコード生成。単純にN個の式を順に生成
        String generate(int indentLevel) {
            return exprs.stream()
                .map(e -> e.generate(indentLevel))
                .collect(Collectors.joining("\n")) + "\n";
        }
    }

    // 順序付き選択
    public static class Choice extends Expr {
        public final Expr alt1, alt2;
        public Choice(Expr alt1, Expr alt2) {
            this.alt1 = alt1;
            this.alt2 = alt2;
        }

        // 順序付き選択のコード生成
        // try-catchを使ってバックトラックを実装
        String generate(int indentLevel) {
            var code = """
                int saved = pos;
                try {
                    %s
                } catch (Failure e) {
                    pos = saved;
                    %s
                }
             """.formatted(
                alt1.generate(indentLevel + 4).stripTrailing(), 
                alt2.generate(indentLevel + 4));
            return code.indent(indentLevel);
        }
    }
    
    // コード生成のメインメソッド
    public static String generateParser(List<Rule> rules) {
        StringBuilder code = new StringBuilder();
        code.append("""
            public class Parser {
                public static class Failure extends RuntimeException {
                    public Failure(String message) {
                        super(message);
                    }
                }
                private String input;
                private int pos;
                public Parser(String input) {
                    this.input = input;
                    this.pos = 0;
                }
                public void match(String str) {
                    if (pos + str.length() <= input.length() &&
                        input.startsWith(str, pos)) {
                        pos += str.length();
                        return;
                    }
                    throw new Failure("Expected '" + str + "' at pos " + pos);
                }
        """.stripLeading());
        // 各規則に対してメソッドを生成
        for (Rule rule : rules) {
            code.append("    public void parse" + rule.name + "() {\n");
            code.append(rule.body.generate(8));
            code.append("\n");
            code.append("        return;\n");
            code.append("    }\n\n");
        }
        code.append("}\n");
        return code.toString();
    }
}
\end{codescreen}

先ほど作ったDyck言語のPEGをちょっと変形したものをこの構文解析器生成系に与えます。利用コードは次のようになります。

\begin{codescreen}{}
// D <- P;
// P <- "(" P ")" P / "";
List<Rule> rules = Arrays.asList(
    new Rule("D", new NT("P")),
    new Rule("P", 
        new Choice(
            new Seq(
                new Lit("("),
                new NT("P"),
                new Lit(")"),
                new NT("P")
            ),
            // 空文字列 ε（イプシロン）
            new Lit("")
         )
    )
);
// 構文解析器のコードを生成
String parserCode = PEG2Java.generateParser(rules);
System.out.println(parserCode);
\end{codescreen}

次のようなJavaコードが生成されます。

\begin{codescreen}{}
public class Parser {
    public static class Failure extends RuntimeException {
        public Failure(String message) {
            super(message);
        }
    }
    private String input;
    private int pos;
    public Parser(String input) {
        this.input = input;
        this.pos = 0;
    }
    public void match(String str) {
        if (pos + str.length() <= input.length() &&
            input.startsWith(str, pos)) {
            pos += str.length();
            return;
        }
        throw new Failure("Expected '" + str + "' at pos " + pos);
    }

    public void parseD() {
        parseP();
        return;
    }

    public void parseP() {
        int saved = pos;
        try {
            match("(");
            parseP();
            match(")");
            parseP();
        } catch (Failure e) {
            pos = saved;
            match("");
        }

        return;
    }
}
\end{codescreen}

空文字列への分岐（\texttt{match("")}）を残しているので、\texttt{()}のように\texttt{P}が空になるケースも、\texttt{(())}のようにネストするケースも成功します。\texttt{try}ブロックで失敗したときに位置を巻き戻してから空文字列を選択する、という5.1.2節と同じ挙動を保っています。

このコードは実際にDyck言語の変形版を解析できます。\texttt{parseD()}メソッドを呼び出すことで、Dyck言語の文字列が正しいかどうかをチェックできます。

このように、PEGの文法定義から機械的にJavaコードを生成できます。本体のコードだけでいえばわずか140行程度です。これでPEGの構文解析器生成系が実装できるのです。もちろん、実用的には抽象構文木（AST）の生成や、エラーメッセージの改善などが必要になりますが、コアの部分はこのようにシンプルに実装できます。

\section{構文解析器生成系の分類}
\index{こうぶんかいせききせいせいけい@構文解析器生成系!分類}
\index{LL(1)}
\index{LALR(1)}
\index{GLR}
\index{じくかいせききせいせいけい@字句解析器生成系}
\index{Lex}
\index{Yacc}

構文解析器生成系は1970年代頃から研究の蓄積があり、数多くの構文解析器生成系がこれまで開発されています。

基本的には、構文解析器生成系は採用している構文解析アルゴリズムによって分類されます。たとえば

\begin{itemize}
\item
  JavaCCはLL(k)構文解析器を出力するため、\textbf{\sffamily LL(k)}\textgt{構文解析器生成系}
\item
  YaccはLALR(1)構文解析器を出力するため、\textbf{\sffamily LALR(1)}\textgt{構文解析器生成系}
\end{itemize}

一般的な構文解析器生成系の処理フローは、おおむね以下のようになります。ただし、PEGを採用している構文解析器生成系では、字句解析器の部分が丸ごと不要になります。

\begin{figure}[htb]
\centering
\begin{tikzpicture}[
    scale=0.7, transform shape,
    node distance=2.2cm,
    box/.style={draw, rectangle, minimum width=3cm, minimum height=0.8cm, align=center, font=\small},
    diamond/.style={draw, ellipse, minimum width=2.8cm, minimum height=1cm, align=center, font=\small},
    arrow/.style={->, thick},
    label/.style={midway, above, font=\footnotesize}
]
    % ノードの定義
    \node[box] (A) {文法定義ファイル\\(.g, .y, .jjなど)};
    \node[diamond, right=1.8cm of A] (B) {構文解析器\\生成系};
    \node[box, above right=1cm and 1.5cm of B] (C) {字句解析器\\コード生成部};
    \node[box, right=1.8cm of C] (D) {生成された\\字句解析器コード\\(.java, .cなど)};
    \node[box, below right=1cm and 1.5cm of B] (E) {構文解析器\\コード生成部};
    \node[box, right=1.8cm of E] (F) {生成された\\構文解析器コード\\(.java, .cなど)};
    \node[diamond, right=3.5cm of B] (H) {コンパイラ};
    \node[box, right=1.8cm of H] (I) {実行可能な\\パーサー};

    % 矢印の定義
    \draw[arrow] (A) -- (B);
    \draw[arrow] (B) -- node[label,sloped] {字句解析ルール} (C);
    \draw[arrow] (C) -- (D);
    \draw[arrow] (B) -- node[label,sloped] {構文解析ルール} (E);
    \draw[arrow] (E) -- (F);
    \draw[arrow] (D) -- (H);
    \draw[arrow] (F) -- (H);
    \draw[arrow] (H) -- (I);
\end{tikzpicture}
\caption{一般的な構文解析器生成系の処理フロー}
\label{fig:parser-generator-flow}
\end{figure}

Yacc/Lexのように、字句解析器生成系（Lex）と構文解析器生成系（Yacc）が別々のツールとして提供され、連携して動作するケースもあります。

ただし、GNU BisonはYaccと違って、LALR(1)より広いGLR（Generalized
LR：一般化LR）構文解析器も生成できるので、GLR構文解析器生成系であるとも言えます。GLRは、通常のLR構文解析ではコンフリクト（シフト/還元コンフリクトなど）が発生するような曖昧な文法も扱えるように拡張された手法です。Bisonを使う場合、ほとんどはLALR(1)構文解析器を出力するので、GLRについて言及されることは少ないですが、知っておいても損はないでしょう。

より大きなくくりでみると、下向き構文解析（LL法やPEG）と上向き構文解析（LR法など）という観点から分類することもできますし、ともに文脈自由文法ベースであるLL法やLR法と、PEGなど他の形式言語を用いた構文解析法を対比してみせることもできます。

以下に、本書で紹介する代表的な構文解析器生成系の比較をまとめます。

\begin{table}[htb]
  \small
  \begin{center}
    \begin{tabular}{p{3cm}p{3cm}p{3cm}p{3cm}}\hline
      特徴項目 & JavaCC & Yacc & ANTLR\\ \hline
      採用アルゴリズム & LL(k)(デフォルトはLL(1)) & LALR(1) & ALL(*)\\ \hline
      生成コードの言語 & Java & C, C++ & Java, C++, Python, JavaScript, Go, C\#, Swift, Dart, PHP (多言語対応)\\ \hline
      左再帰の扱い & 不可(文法書き換えが必要) & 左再帰を扱える(直接・間接とも) & 直接左再帰を扱える(v4以降。間接左再帰は不可)\\ \hline
      曖昧性解決 & 先読みトークン数(k)の調整、意味アクション & 演算子の優先順位・結合規則指定、\%precなどで対応	 & 意味アクション、構文述語、ALL(*)による自動解決\\ \hline
      エラー報告/リカバリ機能 & 基本的 & errorトークンによる限定的なリカバリ & 高度なエラー報告、柔軟なエラーリカバリ戦略\\ \hline
      学習コスト & Javaユーザーには比較的容易 & やや高め(C言語と連携の知識も必要) & 機能が豊富で強力な分、やや高め\\ \hline
      その他特徴 & Javaに特化、構文がJavaライク & C言語との親和性が高い、歴史と実績がある & 強力な解析能力、豊富なターゲット言語、優れたツールサポート(GUIなど)\\ \hline
    \end{tabular}
  \end{center}
\end{table}

\section{JavaCC：Javaの構文解析器生成系の定番}
\index{JavaCC}
\index{LL(1)}
\index{LL(k)}
\index{LOOKAHEAD}
\index{TOKEN}
\index{SKIP}
\index{EOF}
\index{せまんてぃっくあくしょん@セマンティックアクション}

1996年、Sun
Microsystems（当時）は、Jackという構文解析器生成系をリリースしました。その後、Jackの作者が自らの会社を立ち上げ、JackはJavaCCに改名されて広く知られることとなりましたが、現在では紆余曲折の末、\href{https://javacc.github.io/javacc}{javacc.github.io}の元で開発およびメンテナンスが行われています。現在のライセンスは3条項BSDライセンスです。

JavaCCという名前は「Java Compiler
Compiler」の略で、Yacc同様「コンパイラを作るためのコンパイラ」という意味です。

JavaCCはLL(k)法を元に作られており、構文定義ファイルからLL(k)パーサーを生成します。以下は四則演算を含む数式を計算できる電卓をJavaCCで書いた場合の例です。

\begin{codescreen}{}
options {
  STATIC = false;
  JDK_VERSION = "21";
}

PARSER_BEGIN(Calculator)
package com.github.kmizu.calculator;
public class Calculator {
  public static void main(String[] args) throws ParseException {
    Calculator parser = new Calculator(System.in);
    System.out.println(parser.start());
  }
}
PARSER_END(Calculator)

SKIP : { " " | "\t"  | "\r" | "\n" }
TOKEN : {
  <ADD: "+">
| <SUBTRACT: "-">
| <MULTIPLY: "*">
| <DIVIDE: "/">
| <LPAREN: "(">
| <RPAREN: ")">
| <INTEGER: (["0"-"9"])+>
}

public int start() :
{int r = 0;}
{
  r=add() { return r; }
}

public int expression() :
{int r = 0;}
{
  r=add() <EOF> { return r; }
}

public int add() :
{int r = 0; int v = 0;}
{
    r=mult() ( <ADD> v=mult() { r += v; }| <SUBTRACT> v=mult() { r -= v; })* {
        return r;
    }
}


public int mult() :
{int r = 0; int v = 0;}
{
   r=primary() ( 
     <MULTIPLY> v=primary() { r *= v; }
   | <DIVIDE> v=primary() { r /= v; })* { return r; }
}

public int primary() :
{int r = 0; Token t = null;}
{
(
  <LPAREN> r=add() <RPAREN>
| t=<INTEGER> { r = Integer.parseInt(t.image); }
) { return r; }
}
\end{codescreen}

この\texttt{SKIP}から\texttt{TOKEN}までの部分はトークン定義になります。トークンとは、構文解析の基本単位となる要素で、「+」「123」「(」などのことです。第2章のJSONパーサーでも「文字列」「数値」「中括弧」などをトークンとして扱いましたね。ここでは、7つのトークンを定義しています。トークン定義の後が構文規則の定義になります。ここでは、

\begin{itemize}
\item
  \texttt{start()}
\item
  \texttt{expression()}
\item
  \texttt{add()}
\item
  \texttt{mult()}
\item
  \texttt{primary()}
\end{itemize}

\noindent の5つの構文規則が定義されています。各構文規則はJavaのメソッドに酷似した形で記述されますが、実際、ここから生成される.javaファイルには同じ名前のメソッドが定義されます。

\texttt{start()}と\texttt{expression()}は両方とも\texttt{add()}を呼び出すエントリポイントですが、違いは入力末尾の扱い方です。\texttt{expression()}は最後に\texttt{\textless{}EOF\textgreater{}}（入力終端トークン）を要求するため、入力全体が1つの式であることを厳密に検証したい場面（テストやファイル全体の構文チェックなど）に向きます。一方、\texttt{start()}は\texttt{\textless{}EOF\textgreater{}}を要求しないので、標準入力のような対話的な入力で「式を1つ読み終えたら抜ける」という使い方ができます。\texttt{main}メソッドでは標準入力を読むので\texttt{start()}を呼んでいます。

\texttt{add()}が\texttt{mult()}を呼び出して、\texttt{mult()}が\texttt{primary()}を呼び出すという構図は第1章ですでにみた形ですが、第1章と違って単純に宣言的に各構文規則の関係を書けばそれでOKなのが構文解析器生成系の強みです。

ちなみに、\texttt{\{int\ r\ =\ 0;\}}のような部分はJavaコードの埋め込みで、構文解析中に実行される処理を記述しています。これを「セマンティックアクション」（意味的動作）と呼びます。セマンティックアクションを使うことで、単に文法が正しいかチェックするだけでなく、解析しながら計算や抽象構文木の構築などの処理を行えます。

このようにして定義した電卓プログラムは次のようにして利用できます。

\begin{codescreen}{}
package com.github.kmizu.calculator;

import org.junit.jupiter.api.DisplayName;
import org.junit.jupiter.api.Test;
import com.github.kmizu.calculator.Calculator;
import java.io.*;

import static org.junit.jupiter.api.Assertions.assertEquals;

public class CalculatorTest {
    @Test
    @DisplayName("1 + 2 * 3 = 7")
    public void test1() throws Exception {
        Calculator calculator = new Calculator(
            new StringReader("1 + 2 * 3"))
        ;
        assertEquals(7, calculator.expression());
    }

    @Test
    @DisplayName("(1 + 2) * 4 = 12")
    public void test2() throws Exception {
        Calculator calculator = new Calculator(
            new StringReader("(1 + 2) * 4")
        );
        assertEquals(12, calculator.expression());
    }

    @Test
    @DisplayName("(5 * 6) - (3 + 4) = 23")
    public void test3() throws Exception {
        Calculator calculator = new Calculator(
            new StringReader("(5 * 6) - (3 + 4)")
        );
        assertEquals(23, calculator.expression());
    }
}
\end{codescreen}

この\texttt{CalculatorTest}クラスではJUnit5を使って、JavaCCで定義した\texttt{Calculator}クラスの挙動をテストしています。空白や括弧を含む数式を問題なく計算できているのがわかるでしょう。

このようなケースでは先読みトークン数が1のため、JavaCCのデフォルトで構いませんが、定義したい構文によっては先読み数を2以上に増やさなければいけないこともあります。以下のようにして先読み数を増やせます。

\begin{codescreen}{}
options {
  STATIC = false;
  JDK_VERSION = "21";
  LOOKAHEAD = 2;
}
\end{codescreen}

\texttt{LOOKAHEAD\ =\ 2}というオプションによって、先読みトークン数を2に増やしています。

先読みトークン数が2ということは、「次の2つのトークンを見て、どの規則を適用するか決定できる」という意味です。LOOKAHEADは固定されていれば任意の正の整数にできるので、JavaCCはデフォルトではLL(1)だが、オプションを設定することによってLL(k)になるともいえます。

JavaCCは構文定義ファイルの文法がかなりJavaに似ているため、生成されるコードの形を想像しやすいというメリットがあります。Javaの構文解析器生成系の中では最古の部類の割に今でも現役で使われているのは、Javaユーザーにとっての使いやすさが背景にあるように思います。

\section{Yacc：構文解析器生成系の老舗}
\index{Yacc}
\index{Bison}
\index{Lex}
\index{Flex}
\index{LALR(1)}
\index{error token@\texttt{error}トークン}

Yaccは1970年代にAT\&Tのベル研究所にいたStephen C.
Johnson（スティーブン・ジョンソン）によって作られたソフトウェアで、非常に歴史がある構文解析器生成系です。第4章で学んだLALR(1)法の構文解析器生成系を実用化した最初のツールでもあります。

現在広く使われているGNU
BisonはYaccのGNUによる再実装で、GNUプロジェクトの一部として開発されています。GNU
BisonはYaccと互換性があり、Yaccの機能を拡張したものです。

Yaccを使って、四則演算を行う電卓プログラムを作るにはまず字句解析器生成系であるLex（正確には、Vern
Paxsonによる互換実装であるFlex）用の定義ファイルを書く必要があります。Lexの定義ファイル\texttt{token.l}は次のようになります。

\begin{codescreen}{}
  %{
  #include "y.tab.h"
  extern int yylval;
  %}

  %%
  [0-9]+      { yylval = atoi(yytext); return NUM; }
  [-+*/()]    { return yytext[0]; }
  [ \t]       { /* ignore whitespace */ }
  "\r\n"      { return EOL; }
  "\r"        { return EOL; }
  "\n"        { return EOL; }
  .           { printf("Invalid character: %s\n", yytext); }
  %%
\end{codescreen}

\texttt{\%\%}から\texttt{\%\%}までがトークンの定義です。これはそれぞれ次のように読むことができます。

\begin{itemize}
\item
  \texttt{{[}0-9{]}+\ \ \ \ \ \ \{\ yylval\ =\ atoi(yytext);\ return\ NUM;\ \}}

  0-9までの数字が1個以上あった場合は数値として解釈し（\texttt{atoi(yytext)})、トークン\texttt{NUM}として返します。
\item
  \texttt{{[}-+*/(){]}\ \ \ \ \{\ return\ yytext{[}0{]};\ \}}

  \texttt{-},\texttt{+},\texttt{*},\texttt{/},\texttt{(},\texttt{)}についてはそれぞれの文字をそのままトークンとして返します。
\item
  \texttt{{[}\ \textbackslash{}t{]}\ \ \ \ \ \ \ \{\ /*\ ignore\ whitespace\ */\ \}}

  タブおよびスペースは単純に無視します
\item
  \texttt{"\textbackslash{}r\textbackslash{}n"\ \ \ \ \ \ \{\ return\ EOL;\ \}}

  改行はEOLというトークンとして返します
\item
  \texttt{"\textbackslash{}r"\ \ \ \ \ \ \ \ \{\ return\ EOL;\ \}}

  前に同じ
\item
  \texttt{"\textbackslash{}n"\ \ \ \ \ \ \ \ \{\ return\ EOL;\ \}}

  前に同じ
\item
  \texttt{.\ \ \ \ \ \ \ \ \ \ \ \{\ printf("Invalid\ character:\ \%s\textbackslash{}n",\ yytext);\ \}}

  それ以外の文字がきたらエラーを吐きます
\end{itemize}

このトークン定義ファイルを入力として与えると、flexは\texttt{lex.yy.c}という形で字句解析器を出力します。

次に、yaccの構文定義ファイルを書きます。

\begin{codescreen}{}
  %{
  #include <stdio.h>
  int yylex(void);
  void yyerror(char const *s);
  int yywrap(void) {return 1;}
  extern int yylval;
  %}

  %token NUM
  %token EOL
  %left '+' '-'
  %left '*' '/'

  %%

  input :
        | input line
        ;

  line : expr EOL
        { printf("Result: %d\n", $1); }
        | EOL
        | error EOL
        { yyerrok; }
        ;

  expr : NUM
        { $$ = $1; }
        | expr '+' expr
        { $$ = $1 + $3; }
        | expr '-' expr
        { $$ = $1 - $3; }
        | expr '*' expr
        { $$ = $1 * $3; }
        | expr '/' expr
        {
          if ($3 == 0) { yyerror("Cannot divide by zero."); YYERROR; }
          else { $$ = $1 / $3; }
        }
        | '(' expr ')'
        { $$ = $2; }
        ;

  %%

  void yyerror(char const *s)
  {
    fprintf(stderr, "Parse error: %s\n", s);
  }

  int main()
  {
    yyparse();
  }
\end{codescreen}

flexの場合と同じく、\texttt{\%\%}から\texttt{\%\%}までが構文規則の定義の本体です。実行されるコードが入っているので読みづらくなっていますが、それを除くと以下のようになります。

\begin{codescreen}{}
  %token NUM
  %token EOL
  %left '+' '-'
  %left '*' '/'

  %%
  input :
        | input line
        ;

  line  : expr EOL
        | EOL
        | error EOL
        ;

  expr  : NUM
        | expr '+' expr
        | expr '-' expr
        | expr '*' expr
        | expr '/' expr
        | '(' expr ')'
        ;
  %%
\end{codescreen}

だいぶ見やすくなりましたね。入力を表す\texttt{input}規則は、0個以上の\texttt{line}からなります。\texttt{line}は式を表す\texttt{expr}と改行からなるので、式を1行ずつ読み取って結果を出力できます。

次に、\texttt{expr}規則ですが、ここではyaccの優先順位を表現するための機能である\texttt{\%left}を使ったため、優先順位のためだけに規則を作る必要がなくなっており、定義が簡潔になっています。ともあれ、こうして定義された文法定義ファイルをyaccに与えると\texttt{y.tab.c}というファイルを出力します。

flexとyaccが出力したファイルを結合して実行ファイルを作ると、次のような入力に対して

\begin{codescreen}{}
1 + 2 * 3
2 + 3
6 / 2
3 / 0
\end{codescreen}

それぞれ、次のような出力を返します。

\begin{codescreen}{}
Result: 7
Result: 5
Result: 3
Parse error: Cannot divide by zero.
\end{codescreen}

Yaccはとても古いソフトウェアの1つですが、Rubyの文法定義ファイルparse.yはYacc用ですし、未だに各種言語処理系では現役で使われてもいます。

\section{ANTLR：多言語対応の強力な下向き構文解析器生成系}
\index{ANTLR}
\index{PCCTS}
\index{LL(k)}
\index{LL(*)}
\index{ALL(*)}
\index{Adaptive LL(*)}
\index{ひだりさいき@左再帰}
\index{じゅつご@述語}
\index{semantic predicate}

1989年にPurdue Compiler Construction Tool
Set（PCCTS）というものがありました。ANTLRはその後継というべきもので、これまでに、\texttt{LL(k) -> LL(*) -> ALL(*)}と構文解析アルゴリズムを拡張し、取り扱える文法の幅を広げつつアクティブに開発が続けられています。作者はTerence
Parrという方ですが、構文解析器一筋と言っていいくらいに、ANTLRにキャリアの多くを費やしてきている人です。

それだけに、ANTLRの完成度は非常に高いものになっています。また、一時期はLR法に比べてLL法の評価は低いものでしたが、Terence
ParrがLL(k)を改良していく過程で、LL(\texttt{*})やALL(\texttt{*})のようなLR法に比べてもなんら劣らない、実用的にも使いやすい構文解析法が発明されました。

ANTLRはJava、C++などいくつもの言語を扱うことができますが、特に安心して使えるのはJavaです。以下は先ほどと同様の、四則演算を解析できる数式パーサーをANTLRで書いた場合の例です。

ANTLRでは構文規則は、\texttt{規則名\ :\ 本体\ ;}
という形で記述しますが、LLパーサー向けの構文定義を素直に書き下すだけでOKです。

\begin{codescreen}{}
grammar Expression;

expression returns [int e]
    : v=additive {$e = $v.e;}
    ;

additive returns [int e = 0;]
    : l=multitive {$e = $l.e;} (
      '+' r=multitive {$e = $e + $r.e;}
    | '-' r=multitive {$e = $e - $r.e;}
    )*
    ;

multitive returns [int e = 0;]
    : l=primary {$e = $l.e;} (
      '*' r=primary {$e = $e * $r.e;}
    | '/' r=primary {$e = $e / $r.e;}
    )*
    ;

primary returns [int e]
    : n=NUMBER {$e = Integer.parseInt($n.getText());}
    | '(' x=expression ')' {$e = $x.e;}
    ;

LP : '(' ;
RP : ')' ;
NUMBER : INT ;
fragment INT :   '0' | [1-9] [0-9]* ; // no leading zeros
WS  :   [ \t\n\r]+ -> skip ;
\end{codescreen}

規則\texttt{expression}が数式を表す規則です。そのあとに続く\texttt{returns\ {[}int\ e{]}}はこの規則を使って解析を行った場合に\texttt{int}型の値を返すことを意味しています。これまで見てきたように構文解析をした後には抽象構文木をはじめとして何らかのデータ構造を返す必要があります。\texttt{returns\ ...}はそのために用意されている構文です。名前がすべて大文字の規則はトークンを表しています。

数式を表す各規則についてはこれまで書いてきた構文解析器と同じ構造なので読むのに苦労はしないでしょう。

規則\texttt{WS}は空白文字を表すトークンですが、これは数式を解析する上では読み飛ばす必要があります。\\
\texttt{{[}\ \textbackslash{}t\textbackslash{}n\textbackslash{}r{]}+\ -\textgreater{}\ skip}は

\begin{itemize}
\item
  スペース
\item
  タブ文字
\item
  改行文字
\end{itemize}

\noindent のいずれかが出現した場合は読み飛ばすということを表現しています。

ANTLRは下向き型の構文解析が苦手とする左再帰も、規則の中で直接的に自分自身を呼び出す形（直接左再帰）であれば扱えます（規則を経由した間接左再帰には対応していない点には注意してください）。先ほどの定義ファイルでは繰り返しを使っていましたが、これを直接左再帰に直した以下の定義ファイルも全く同じ挙動をします。

\begin{codescreen}{}
grammar LRExpression;

expression returns [int e]
    : v=additive {$e = $v.e;}
    ;

additive returns [int e]
    : l=additive op='+' r=multitive {$e = $l.e + $r.e;}
    | l=additive op='-' r=multitive {$e = $l.e - $r.e;}
    | v=multitive {$e = $v.e;}
    ;

multitive returns [int e]
    : l=multitive op='*' r=primary {$e = $l.e * $r.e;}
    | l=multitive op='/' r=primary {$e = $l.e / $r.e;}
    | v=primary {$e = $v.e;}
    ;

primary returns [int e]
    : n=NUMBER {$e = Integer.parseInt($n.getText());}
    | '(' x=expression ')' {$e = $x.e;}
    ;

LP : '(' ;
RP : ')' ;
NUMBER : INT ;
fragment INT :   '0' | [1-9] [0-9]* ; // no leading zeros
WS  :   [ \t\n\r]+ -> skip ;
\end{codescreen}

左再帰を使うことでより簡単に文法を定義できることもあるので、あるとうれしい機能だと言えます。

さらに、ANTLRは\texttt{ALL(*)}というアルゴリズムを採用しているため、通常のLLパーサーでは扱えないような文法定義も取り扱えます。以下の「最小XML」文法定義を見てみましょう。

\begin{codescreen}{}
grammar PetitXML;
@parser::header {
    import static com.github.asciidwango.parser_book.ch5.PetitXML.*;
    import java.util.*;
}

root returns [Element e]
    : v=element {$e = $v.e;}
    ;

element returns [Element e]
    : ('<' begin=NAME '>' es=elements '</' end=NAME '>' {
          $begin.text.equals($end.text)
      }?
      {$e = new Element($begin.text, $es.es);})
    | ('<' name=NAME '/>' {$e = new Element($name.text);})
    ;

elements returns [List<Element> es]
    : { $es = new ArrayList<>();} (element {$es.add($element.e);})*
    ;

LT:    '<';
GT:    '>';
SLASH: '/';
NAME:  [a-zA-Z_][a-zA-Z0-9]* ;

WS :   [ \t\n\r]+ -> skip ;
\end{codescreen}

\texttt{PetitXML}の名のとおり、属性やテキストなどは全く扱えず、\texttt{\textless{}a\textgreater{}}や\texttt{\textless{}a/\textgreater{}}、\texttt{\textless{}a\textgreater{}\textless{}b\textgreater{}\textless{}/b\textgreater{}\textless{}/a\textgreater{}}といった要素のみを扱えます。規則\texttt{element}が重要です。

\begin{codescreen}{}
element returns [Element e]
    : ('<' begin=NAME '>' es=elements '</' end=NAME '>' {
            $begin.text.equals($end.text)
       }?
      {$e = new Element($begin.text, $es.es);})
    | ('<' name=NAME '/>' {$e = new Element($name.text);})
    ;
\end{codescreen}

ここで空要素（\texttt{\textless{}e/\textgreater{}}など）に分岐するか、子要素を持つ要素（\texttt{\textless{}a\textgreater{}\textless{}b/\textgreater{}\textless{}/a\textgreater{}}など）に分岐するかを決定するには、\texttt{\textless{}}に加えて、任意の長さになり得るタグ名を先読みしなければいけません。通常のLLパーサーでは何文字（何トークン）先読みできるかは予め決定されているのでこのような文法定義を取り扱うことはできません。しかし、ANTLRの\texttt{ALL(*)}アルゴリズムやその前身となる\texttt{LL(*)}アルゴリズムでは任意個の文字数を先読みして分岐を決定できます。

ANTLRでは通常のLLパーサーと違い文法を記述する上での大きな制約がなく、非常に強力です。理論的な意味でも\texttt{ALL(*)}アルゴリズムは任意の決定的な文脈自由言語を取り扱えます。

また、\texttt{ALL(*)}アルゴリズム自体とは関係ありませんが、XMLのパーサーを書くときには開きタグと閉じタグの名前が一致している必要があります。この条件を記述するために\texttt{PetitXML}では次のように記述されています。

\begin{codescreen}{}
'<' begin=NAME '>' es=elements '</' end=NAME '>' {
  $begin.text.equals($end.text)
}?
\end{codescreen}

この中の\texttt{\{\$begin.text.equals(\$end.text)\}?}という部分はsemantic
predicate（意味的述語）と呼ばれ、プログラムとして書かれた条件式が真になるときにだけマッチします。この例では、開きタグの名前（\texttt{begin}）と閉じタグの名前（\texttt{end}）が一致しているかをJavaのコードで確認しています。semantic
predicateのような機能はプログラミング言語をそのまま埋め込むという意味で、正直「あまりきれいではない」と思わなくもないですが、実用上はsemantic
predicateを使いたくなる場面にしばしば遭遇します。

ANTLRはこういった実用上重要な痒いところにも手が届くように作られており、非常によくできた構文解析器生成系といえるでしょう。

\section{SComb：Scala製のパーサーコンビネータ}
\index{SComb}
\index{ぱーさーこんびねーた@パーサーコンビネータ}
\index{Scala}
\index{DSL}

手前味噌ですが、拙作のパーサーコンビネータである\href{https://github.com/kmizu/scomb}{SComb}も紹介しておきます。これまで紹介してきたものはすべて構文解析器生成系です。つまり、独自の言語を用いて作りたい言語の文法を記述し、そこから\textgt{対象言語}（CであったりJavaであったりさまざまですが）で書かれた構文解析器を生成するものだったわけですが、パーサーコンビネータは少々違います。

パーサーコンビネータでは対象言語のメソッドや関数、オブジェクトとして構文解析器を定義し、演算子やメソッドによって構文解析器を組み合わせることで構文解析器を組み立てていきます。「コンビネータ」という名前は「組み合わせる（combine）」からきており、小さなパーサーを組み合わせて大きなパーサーを作るという意味です。パーサーコンビネータではメソッドや関数、オブジェクトとして規則自体を記述するため、特別にプラグインを作らなくてもIDEによる支援が受けられることや、対象言語が静的型システムを持っていた場合、型チェックによる支援を受けられることなどがメリットとして挙げられます。

SCombで四則演算を解析できるプログラムを書くと以下のようになります。先ほど述べたようにSCombはパーサーコンビネータであり、これ自体がScalaのプログラム（\texttt{object}宣言）でもあります。

\begin{codescreen}{}
object Calculator extends SCombinator {
  // root <- E 
  def root: P[Int] = E

  // E <- A
  def E: P[Int] = rule(A)

  // A <- M ("+" M / "-" M)* 
  def A: P[Int] = rule(chainl(M) {
    $("+").map { op => (lhs: Int, rhs: Int) => lhs + rhs } |
    $("-").map { op => (lhs: Int, rhs: Int) => lhs - rhs }
  })

  // M <- P ("*" P / "/" P)* 
  def M: P[Int] = rule(chainl(P) {
    $("*").map { op => (lhs: Int, rhs: Int) => lhs * rhs } |
    $("/").map { op => (lhs: Int, rhs: Int) => lhs / rhs }
  })

  // P <- "(" E ")" | N
  def P: P[Int] = rule{
    (for {
      _ <- string("("); e <- E; _ <- string(")")} yield e) | N
  }
  
  // N <- [0-9]+ 
  def N: P[Int] = rule(set('0'to'9').+.map{ digits => digits.mkString.toInt})

  def parse(input: String): Result[Int] = parse(root, input)
}
\end{codescreen}

ScalaはJavaとは異なるプログラミング言語ですが、構文にJavaと似た部分が多く（JVM上で動く処理系がある）、Javaユーザーの方でも比較的読み進めやすいはずです。

各メソッドに対応するBNFによる規則をコメントとして付加してみましたが、BNFと比較しても簡潔に記述できています。Scalaは記号をそのままメソッドとして記述できるなど、元々DSL（ドメイン特化言語）に向いている特徴を持った言語なのですが、その特徴を活用しています。\texttt{chainl}というメソッドについてだけは見慣れない読者の方が多そうですが、これは

\begin{codescreen}{}
// A ::= M ("+" M | "-" M)*
\end{codescreen}

\noindent のような二項演算を簡潔に記述するためのコンビネータ（メソッド）です。パーサーコンビネータの別のメリットとして、BNF（あるいはPEG）に無い演算子を後付けで導入できることも挙げられます。構文規則からの値（意味値）の取り出しもScalaのfor式（Javaのfor文とは異なり、値を生成する式です）を用いて簡潔に記述できています。

筆者は自作言語Klassicの処理系作成のためにSCombを使っていましたが、かなり複雑な文法を記述できるにもかかわらず、SCombのコア部分はわずか600行ほどです。それでいて高い拡張性や簡潔な記述が可能なのは、Scalaという言語の能力と、SCombがベースとして利用しているPEGという手法のシンプルさがあってのものだと言えるでしょう。

\section{パーサーコンビネータJCombを自作しよう！}
\index{JComb}
\index{ぱーさーこんびねーた@パーサーコンビネータ!JComb}
\index{じぇねりくす@ジェネリクス}

コンパイラについて解説した本は数えきれないほどありますし、その中で構文解析アルゴリズムについて説明した本も少なからずあります。しかし、構文解析アルゴリズムについてのみフォーカスした本はParsing Techniquesほぼ一冊といえる現状です。その上でパーサーコンビネータの自作まで踏み込んだ書籍はほぼ皆無と言っていいでしょう。読者の方には「さすがにちょっとパーサーコンビネータの自作は無理があるのでは」と思われた方もいるのではないでしょうか。

しかし、現代的な言語であればパーサーコンビネータを自作するのは本当に簡単です。多くの読者の方々が拍子抜けしてしまうくらいに。この節ではJavaで書かれたパーサーコンビネータJCombを自作する過程を通じて皆さんにパーサーコンビネータとはどのようなものかを学んでいただきます。

パーサーコンビネータと構文解析器生成系の関係は、次のように考えられます。

\begin{itemize}
\item
  構文解析器生成系：文法定義ファイル → ツール → Javaコード
\item
  パーサーコンビネータ：Javaコードで直接文法を表現
\end{itemize}

つまり、パーサーコンビネータは「生成」のステップを省き、プログラミング言語の機能を活用して直接パーサーを組み立てる手法なのです。

まず復習になりますが、構文解析器というのは文字列を入力として受け取って、解析結果を返す関数（あるいはオブジェクト）とみなせるのでした。これはパーサーコンビネータ、特にPEGを使ったパーサーコンビネータを実装するときに有用な見方です。この「構文解析器はオブジェクトである」を文字通りとって、以下のようなジェネリックなインタフェース\texttt{JParser\textless{}R\textgreater{}}を定義します。

\begin{codescreen}{}
interface JParser<R> {
  Result<R> parse(String input);
}
\end{codescreen}

ここで構文解析器を表現するインタフェース\texttt{JParser\textless{}R\textgreater{}}は型パラメータ\texttt{R}を受け取ることに注意してください。型パラメータ\texttt{R}は、「このインタフェースは何かの型\texttt{R}に対して動作しますが、具体的な型は使用時に決まります」という意味です。

構文解析の結果は抽象構文木になりますが、インタフェースを定義する時点では抽象構文木がどのような形になるかはわかりようがないので、型パラメータにしておくのです。\texttt{JParser\textless{}R\textgreater{}}はたった1つのメソッド\texttt{parse()}を持ちます。\texttt{parse()}は入力文字列\texttt{input}を受け取り、解析結果を\texttt{Result\textless{}R\textgreater{}}として返します。

\texttt{JParser\textless{}R\textgreater{}}の実装はいったいぜんたいどのようなものになるの？という疑問を脇に置いておけば理解は難しくないでしょう。次に解析結果\texttt{Result\textless{}V\textgreater{}}をレコードとして定義します。

\begin{codescreen}{}
record Result<V>(V value, String rest){}
\end{codescreen}

レコード\texttt{Result\textless{}V\textgreater{}}は解析結果を保持するクラスです。\texttt{value}は解析結果の値を表現し、\texttt{rest}は解析した結果「残った」文字列を表します。たとえば、\texttt{"123abc"}という文字列から数値部分\texttt{"123"}を解析した場合、\texttt{value}は\texttt{123}（整数値）、\texttt{rest}は\texttt{"abc"}（残りの文字列）となります。

このインタフェース\texttt{JParser\textless{}R\textgreater{}}は次のように使えると理想的です。

\begin{codescreen}{}
JParser<Integer> calculator = ...;
Result<Integer> result = calculator.parse("1+2*3");
assert 7 == result.value();
\end{codescreen}

パーサーコンビネータは、このようなどこか都合のよい\texttt{JParser\textless{}R\textgreater{}}を、BNF（あるいはPEG）に近い文法規則を連ねていくのに近い使い勝手で構築するための技法です。前の節で紹介した\texttt{SComb}もパーサーコンビネータでしたが基本的には同じようなものです。

この節では最終的に上のような式を解析できるパーサーコンビネータを作るのが目標です。

\subsection{部品を考えよう}

これからパーサーコンビネータを作っていくわけですが、パーサーコンビネータの基本となる「部品」を作る必要があります。

まず最初に、文字列リテラルを受け取ってそれを解析できる次のような\texttt{string()}メソッドはぜひともほしいところです。

\begin{codescreen}{}
assert new Result<String>("123", "").equals(string("123").parse("123"));
\end{codescreen}

これはBNFで言えば文字列リテラルの表記に相当します。

次に、解析に成功したとしてその値を別の値に変換するための方法もほしいところです。

たとえば、\texttt{123}という文字列を解析したとして、Stringではなくintに変換したいところです。Javaのラムダ式（無名関数）を使えば、このような変換を簡潔に書けます。このようなメソッドは、ラムダ式で変換を定義できるように、次のような\texttt{map()}メソッドとして提供したいところです。

\begin{codescreen}{}
<T, U> JParser<U> map(JParser<T> parser, Function<T, U> function);
assert (new Result<Integer>(123, "")).equals(
    map(string("123"), v -> Integer.parseInt(v)).parse("123")
);
\end{codescreen}

これは構文解析器生成系でセマンティックアクションを書くのに相当すると言えます。

BNFで\texttt{a\ \textbar{}\ b}、つまり選択を書くのに相当するメソッドも必要です。これは次のような\texttt{alt()}メソッドとして提供します。

\begin{codescreen}{}
<T> JParser<T> alt(JParser<T> p1, JParser<T> p2);
assert (new Result<String>("bar", "")).equals(
    alt(string("foo"), string("bar")).parse("bar")
);
\end{codescreen}

同様に、BNFで\texttt{a\ b}、つまり連接を書くのに相当するメソッドも必要ですが、これは次のような\texttt{seq()}メソッドとして提供します。

\begin{codescreen}{}
record Pair<A, B>(A a, B b){}
<A, B> JParser<Pair<A, B>> seq(JParser<A> p1, JParser<B> p2);
assert (new Result<>(new Pair<>("foo", "bar"), "")).equals(
    seq(string("foo"), string("bar")).parse("foobar")
);
\end{codescreen}

最後に、BNFで\texttt{a*}、つまり0回以上の繰り返しに相当する\texttt{rep0()}メソッドや、

\begin{codescreen}{}
<T> JParser<List<T>> rep0(JParser<T> p);
\end{codescreen}

\texttt{a+}、つまり1回以上の繰り返しに相当する\texttt{rep1()}メソッドもほしいところです。

\begin{codescreen}{}
<T> JParser<List<T>> rep1(JParser<T> p);
assert (new Result<List<String>>(List.of("a", "a", "a"), "")).equals(
  rep1(string("a")).parse("aaa")
);
\end{codescreen}

この節ではこれらのプリミティブなメソッドの実装方法について説明していきます。

\subsection{string()メソッド}
\index{JComb}
\index{もじれつりてらる@文字列リテラル}

最初に\texttt{string(String\ literal)}メソッドで返す\texttt{JLiteralParser\textless{}String\textgreater{}}の中身を作ってみましょう。\texttt{JLiteralParser}クラスはただ1つのメソッド\texttt{parse()}を持ち、次のような実装になります。

\begin{codescreen}{}
class JLiteralParser implements JParser<String> {
  private String literal;
  public JLiteralParser(String literal) {
    this.literal = literal;

  }
  public Result<String> parse(String input) {
    if(input.startsWith(literal)) {
      return new Result<String>(literal, input.substring(literal.length()));
    } else {
      return null;
    }
  }
}
\end{codescreen}

このクラスは次のようにして使います。

\begin{codescreen}{}
assert new Result<String>("foo", "").equals(new JLiteralParser("foo").parse("foo"));
\end{codescreen}

リテラルを表すフィールド\texttt{literal}が\texttt{input}の先頭とマッチした場合、\texttt{literal}と残りの文字列からなる\texttt{Result\textless{}String\textgreater{}}を返します。そうでない場合は返すべき\texttt{Result}がないので\texttt{null}を返します。簡単ですね。

\texttt{startsWith}メソッドは、文字列がある文字列で始まるかを判定するJavaの標準メソッドです。\texttt{substring}メソッドは、文字列の一部を切り出すメソッドです。

ここでは簡潔さを優先して失敗時に\texttt{null}を返していますが、\texttt{Success}/\texttt{Failure}のような直和型にしておくと、戻り値に\texttt{null}が紛れ込む心配がなくなります。JCombでは説明を短くするため\texttt{null}にしていますが、実用では直和型を作るのがおすすめです。

あとはこのクラスのインスタンスを返す\texttt{string()}メソッドを作成するだけです。利便性のため、以降では各種メソッドはクラス\texttt{JComb}のstaticメソッドとして実装していきます。

\begin{codescreen}{}
public class JComb {
  public static JParser<String> string(String literal) {
    return new JLiteralParser(literal);
  }
}
\end{codescreen}

使うときは次のようになります。

\begin{codescreen}{}
JParser<String> foo = string("foo");
assert new Result<String>("foo", "_bar").equals(foo.parse("foo_bar"));
assert null == foo.parse("baz");
\end{codescreen}

\subsection{alt()メソッド}
\index{せんたく@選択}
\index{じゅんじょつきせんたく@順序付き選択}

2つのパーサーを取って「選択」パーサーを返すメソッド\texttt{alt()}を実装します。先ほどのようにクラスを実装してもいいですが、メソッドは1つだけなのでラムダ式（Java
8から導入された無名関数）にします。

\begin{codescreen}{}
public class JComb {
    // p1 / p2 (PEGの順序付き選択に対応)
    public static <A> JParser<A> alt(JParser<A> p1, JParser<A> p2) {
        return (input) -> {
            var result = p1.parse(input);//(1) p1を試す
            if(result != null) return result;
                                   //(2) p1が成功したらその結果を返す
            return p2.parse(input);//(3) p1が失敗したらp2を試す
        };
    }
}
\end{codescreen}

ラムダ式について復習しておきます。これは実質的に以下の匿名クラスを書いたのと同じになります。

\begin{codescreen}{}
public class JComb {
    public static <A> JParser<A> alt(JParser<A> p1, JParser<A> p2) {
        return new JAltParser<A>(p1, p2);
    }
}

class JAltParser<A> implements JParser<A> {
  private JParser<A> p1, p2;
  public JAltParser(JParser<A> p1, JParser<A> p2) {
    this.p1 = p1;
    this.p2 = p2;
  }
  public Result<A> parse(String input) {
    var result = p1.parse(input);
    if(result != null) return result;
    return p2.parse(input);
  }
}
\end{codescreen}

この定義では、

\begin{enumerate}
\item
  まずパーサー\texttt{p1}を試す
\item
  \texttt{p1}が成功した場合は\texttt{p2}を試すことなく値をそのまま返す
\item
  \texttt{p1}が失敗した場合、\texttt{p2}を試しその値を返す
\end{enumerate}

\noindent という挙動をします。これはBNFの選択 \texttt{\textbar{}}
とは異なり、PEGの「順序付き選択 \texttt{/}」に対応します。

実は今ここで作っているパーサーコンビネータである\texttt{JComb}は（\texttt{SComb}と同様に）PEGをベースとしたパーサーコンビネータだったのです。

PEGベースでないパーサーコンビネータを作ることもできるのですが実装がかなり複雑になってしまいます。PEGの挙動をそのままプログラミング言語に当てはめるのは非常に簡単であるため、今回はPEGを採用しましたが、もし興味があればBNFベース（文脈自由文法ベース）のパーサーコンビネータも作ってみてください。

\subsection{seq()メソッド}
\index{れんせつ@連接}

次に2つのパーサーを取って「連接」パーサーを返すメソッド\texttt{seq()}を実装します。これはPEGの連接
\texttt{e1\ e2} に対応します。先ほどと同じくラムダ式にしてみます。

\begin{codescreen}{}
record Pair<A, B>(A a, B b){}
// p1 p2 (PEGの連接に対応)
public class JComb {
    public static <A, B> JParser<Pair<A, B>> seq(
        JParser<A> p1, JParser<B> p2) {
        return (input) -> {
            var result1 = p1.parse(input); //(1-1) p1を試す
            if(result1 == null) return null; 
                                           //(1-2) p1が失敗したら全体も失敗
            var result2 = p2.parse(result1.rest()); 
                                           //(2-1) p1の残り入力でp2を試す
            if(result2 == null) return null; 
                                           //(2-2) p2が失敗したら全体も失敗
            // 両方成功したら、結果をペアにして返す
            return new Result<>(
                new Pair<A, B>(
                    result1.value(), result2.value()
                ), result2.rest()
            );//(2-3)
        };
    }
}
\end{codescreen}

先ほどの\texttt{alt()}メソッドと似通った実装ですが、\texttt{p1}が失敗したら全体が失敗するのがポイントです。\texttt{p1}と\texttt{p2}の両方が成功した場合は、2つの値のペアを返しています。

\subsubsection{rep0(), rep1()メソッド}
\index{くりかえし@繰り返し}

残りは0回以上の繰り返し（PEGの
\texttt{p*}）を表す\texttt{rep0()}と1回以上の繰り返し（PEGの
\texttt{p+}）を表す\texttt{rep1()}メソッドです。

まず、\texttt{rep0()}メソッドは次のようになります。

\begin{codescreen}{}
public class JComb {
    public static <A> JParser<List<A>> rep0(JParser<A> p) {
        return (input) -> {
            var result = p.parse(input); // (1)
            if(result == null) return new Result<>(List.of(), input); //(2)
            var value = result.value();
            var rest = result.rest();
            var result2 = rep0(p).parse(rest); //(3)
            if(result2 == null) return new Result<>(List.of(value), rest);
            List<A> values = new ArrayList<>();
            values.add(value);
            values.addAll(result2.value());
            return new Result<>(values, result2.rest());
        };
    }
}
\end{codescreen}

\texttt{(1)}でまずパーサー\texttt{p}を適用しています。ここで失敗した場合、0回の繰り返しにマッチしたことになるので、空リストからなる結果を返します（\texttt{(2)}）。そうでなければ、1回以上の繰り返しにマッチしたことになるので、繰り返し同じ処理をする必要がありますが、これは再帰呼出しによって簡単に実装できます（\texttt{(3)}）。シンプルな実装ですね。

\texttt{rep1(p)}は意味的には\texttt{seq(p,\ rep0(p))}なので、次のようにして実装を簡略化できます。

\begin{codescreen}{}
public class JComb {
    public static <A> JParser<List<A>> rep1(JParser<A> p) {
        JParser<Pair<A, List<A>>> rep1Sugar = seq(p, rep0(p));
        return (input) -> {
            var result = rep1Sugar.parse(input);//(1)
            if(result == null) return null;//(2)
            var pairValue = result.value(); // resultから値を取得
            var values = new ArrayList<A>();
            values.add(pairValue.a()); // Pairの最初の要素
            values.addAll(pairValue.b()); // Pairの2番目の要素 (List)
            return new Result<>(values, result.rest()); //(3)
        };
    }
}
\end{codescreen}

実質的な本体は\texttt{(1)}だけで、あとは結果の値が\texttt{Pair}なのを\texttt{List}に加工しているだけですね。

\subsection{map()メソッド}
\index{せまんてぃっくあくしょん@セマンティックアクション}
\index{らむだしき@ラムダ式}

パーサーを加工して別の値を生成するためのメソッド\texttt{map()}を実装してみましょう。\texttt{map()}は\texttt{JParser\textless{}R\textgreater{}}のメソッドとして実装するとメソッドチェインが使えて便利なので、インタフェースの\texttt{default}メソッドとして実装します。

\begin{codescreen}{}
interface JParser<R> {
    Result<R> parse(String input);

    default <T> JParser<T> map(Function<R, T> f) {
        return (input) -> {
           var result = this.parse(input);
           if (result == null) return null;
           return new Result<>(f.apply(result.value()),result.rest());//(1)
        };
    }
}
\end{codescreen}

\texttt{(1)}で\texttt{f.apply(result.value())}として値を加工しているのがポイントです。

\subsection{lazy()メソッド}
\index{ちえんひょうか@遅延評価}
\index{むげんさいき@無限再帰}
\index{そうごさいき@相互再帰}

パーサーを遅延評価するためのメソッド\texttt{lazy()}も導入します。

\begin{codescreen}{}
public class JComb {
    public static <A> JParser<A> lazy(Supplier<JParser<A>> supplier) {
        return (input) -> supplier.get().parse(input);
    }
}
\end{codescreen}

\texttt{lazy}メソッドの必要性について補足します。Javaは先行評価（関数が呼ばれたらすぐに評価）を行う言語であるため、再帰的な文法規則を直接メソッド呼び出しで表現しようとすると、パーサーオブジェクトの構築時に無限再帰が発生し
\texttt{StackOverflowError} となることがあります。たとえば、算術式の文法で
\texttt{expression} が \texttt{additive} を呼び出し、\texttt{additive}
が \texttt{primary} を呼び出し、\texttt{primary} が括弧表現の中で再び
\texttt{expression} を呼び出すような相互再帰構造を考えてみましょう。

\begin{codescreen}{}
// JComb を使った算術式のパーサー定義（簡略版・lazyなしのイメージ）
// 仮に直接代入しようとすると...
public static JParser<Integer> expression() {
    return additive(); 
}
public static JParser<Integer> additive() {
    return primary();   
}
public static JParser<Integer> primary() {
    return alt(
        number,
        seq(string("("), expression(), string(")"))
    );
}
\end{codescreen}

上記のように単純にメソッド呼び出しでパーサーを組み合わせようとすると、\texttt{expression}
の初期化時に \texttt{additive} が必要になり、その \texttt{additive}
の初期化に \texttt{primary} が、さらにその \texttt{primary} の初期化に
\texttt{expression}
が必要となり、循環参照によって初期化が終わらなくなります。

\texttt{lazy} は
\texttt{Supplier\textless{}JParser\textless{}A\textgreater{}\textgreater{}}
を引数に取ることで、実際の \texttt{JParser\textless{}A\textgreater{}}
オブジェクトの取得（\texttt{supplier.get()}）を、そのパーサーが実際に
\texttt{parse()}
メソッドで使われるときまで遅延させます。\texttt{Supplier}はJavaの関数型インタフェースで、「引数なしで値を返す関数」を表します。これにより、相互再帰するパーサー定義でも、オブジェクト構築時の無限再帰を避けられます。算術式の例では、\texttt{expression}
の定義内で \texttt{additive} を呼び出す部分を
\texttt{lazy(()\ -\textgreater{}\ additive())}
のように記述することで、この問題を解決します。

\subsection{regex()メソッド}
\index{せいきひょうげん@正規表現}

せっかくなので正規表現を扱うメソッド\texttt{regex()}も導入してみましょう。

\begin{codescreen}{}
public class JComb {
    public static JParser<String> regex(String regex) {
        return (input) -> {
            var matcher = Pattern.compile(regex).matcher(input);
            if(matcher.lookingAt()) {
                return new Result<>(
                    matcher.group(), input.substring(matcher.end())
                );
            } else {
                return null;
            }
        };
    }
}
\end{codescreen}

引数として与えられた文字列を\texttt{Pattern.compile()}で正規表現に変換して、マッチングを行うだけです。\texttt{matcher.lookingAt()}は文字列の先頭から正規表現にマッチする部分があるかを確認するメソッドです。これは次のようにして使うことができます。

\begin{codescreen}{}
var number = regex("[0-9]+").map(v -> Integer.parseInt(v));
assert (new Result<Integer>(10, "")).equals(number.parse("10"));
\end{codescreen}

\subsection{算術式のインタプリタを書いてみる}

ここまでで作ったクラス\texttt{JComb}と\texttt{JParser}などを使って簡単な算術式のインタプリタを書いてみましょう。仕様は次のとおりです。

\begin{itemize}
\item
  扱える数値は整数のみ
\item
  演算子は加減乗除（\texttt{+\textbar{}-\textbar{}*\textbar{}/}）のみ
\item
  \texttt{()}によるグルーピングができる
\end{itemize}

実装だけを提示すると次のようになります。

\begin{codescreen}{}
public class Calculator {
    // expression は加減算を担当 (左結合)
    // PEG: expression <- multitive ( ( "+" / "-" ) multitive )*
    public static JParser<Integer> expression() {
        return seq( // multitive と (( "+" / "-" ) multitive )* の連接
            lazy(() -> multitive()), // 左辺の multitive (乗除の項)
            rep0( // 0回以上の繰り返し
                seq( // ( "+" / "-" ) と multitive の連接
                    // "+" または "-"
                    alt(string("+"), string("-")), 
                    lazy(() -> multitive()) // 右辺の multitive
                )
            )
        ).map(p -> { // 解析結果を処理するラムダ式
            var left = p.a(); // 初期値 (最初の項)
            var rights = p.b(); // 残りの演算子と項のペアのリスト
            for (var rightPair : rights) { 
                                // 各 Pair<String, Integer> について
                var op = rightPair.a(); // 演算子 ( "+" または "-" )
                var rightValue = rightPair.b(); // 項の値
                if (op.equals("+")) {
                    left += rightValue;
                } else { // op.equals("-")
                    left -= rightValue;
                }
            }
            return left; // 計算結果
        });
    }

    // multitive は乗除算を担当 (左結合)
    // PEG: multitive <- primary ( ( "*" / "/" ) primary )*
    public static JParser<Integer> multitive() {
        return seq( // primary と ( ( "*" / "/" ) primary )* の連接
            lazy(() -> primary()), // 左辺の primary
            rep0( // 0回以上の繰り返し
                seq( // ( "*" / "/" ) と primary の連接
                    alt(string("*"), string("/")), // "*" または "/"
                    lazy(() -> primary()) // 右辺の primary
                )
            )
        ).map(p -> { // 解析結果を処理するラムダ式
            var left = p.a(); // 初期値 (最初の因子)
            var rights = p.b(); // 残りの演算子と因子のペアのリスト
            for (var rightPair : rights) {
                var op = rightPair.a(); // 演算子 ( "*" または "/" )
                var rightValue = rightPair.b(); // 因子の値
                if (op.equals("*")) {
                    left *= rightValue;
                } else { // op.equals("/")
                    if (rightValue == 0) 
                        throw new ArithmeticException("Division by zero"); 
                    left /= rightValue;
                }
            }
            return left; // 計算結果
        });
    }

    // primary <- number / "(" expression ")"
    public static JParser<Integer> primary() {
        return alt( // number または "(" expression ")" の選択
            number, // 数値パーサー
            seq( // "(" expression の連接
                string("("), // 開き括弧
                // 括弧内の式 (expressionを再帰呼び出し)
                seq(
                    lazy(() -> expression()),
                    string(")") // 閉じ括弧
                )
            ).map(p -> p.b().a()) // ((, (expression, ))) の中央の値を
                                  // 取り出す
        );
    }
    
    // number <- [0-9]+ (PEGの正規表現リテラルに対応)
    private static JParser<Integer> number = 
        regex("[0-9]+").map(Integer::parseInt);
}
\end{codescreen}

表記は冗長なもののほぼPEGに一対一に対応しているのがわかるのではないでしょうか？

これに対してJUnitを使って以下のようなテストコードを記述してみます。無事、意図通りに解釈されていることがわかります。

\begin{codescreen}{}
assertEquals(
    new Result<>(7, ""), Calculator.expression().parse("1+2*3")
); // テストをパス (実際には multitive で処理される)
assertEquals(
    new Result<>(0, ""), Calculator.expression().parse("1+2-3")
); // テストをパス
\end{codescreen}

\subsection{自作パーサーコンビネータのススメ}

DSL（ドメイン特化言語）に向いたScalaに比べれば冗長になったものの、手書きで再帰下降パーサーを組み立てるのに比べると大幅に簡潔な記述を実現できました。しかも、JComb全体を通しても500行にすら満たないのは特筆すべきところです。Javaがユーザー定義の中置演算子（\texttt{+}や\texttt{*}のような演算子を自分で定義できる機能）をサポートしていればもっと簡潔にできたのですが、そこは向き不向きといったところでしょうか。

手書きでパーサーを書く必要に迫られた場合でも、パーサーコンビネータを自作すれば、気軽にパーサーを組み立てるためのDSL（Domain Specific Language）を定義できるのです。それだけでなく、Javaのような静的型付き言語を使った場合、IDEによる支援も受けられますし、BNFやPEGにはない便利な演算子を自分で導入することもできます。

パーサーコンビネータはお手軽なだけあって各種プログラミング言語に実装されています。たとえば、Java用なら\href{https://github.com/jparsec/jparsec}{jparsec}があります。しかし、筆者としては、パーサーコンビネータが動作する仕組みを理解するために、ぜひとも\textgt{自分だけの}パーサーコンビネータを実装してみてほしいと思います。

\section{まとめ}

この章では、構文解析器生成系（パーサージェネレータ）という、文法定義から構文解析器を自動生成するソフトウェアについて学びました。

まず、PEGから構文解析器を生成する基本的な仕組みを、Dyck言語を例に具体的に見ていきました。文法規則を機械的にJavaコードに変換することで、手書きの煩雑さから解放され、文法定義に集中できることを確認しました。

次に、代表的な構文解析器生成系として以下を紹介しました。

\begin{itemize}
\item
  \textbf{\sffamily JavaCC}：\texttt{LL(k)}法を採用し、Javaに特化した構文がJavaユーザーに親しみやすい
\item
  \textbf{\sffamily Yacc}：\texttt{LALR(1)}法を採用する歴史ある構文解析器生成系で、C言語との親和性が高い
\item
  \textbf{\sffamily ANTLR}：\texttt{ALL(*)}アルゴリズムにより強力な解析能力を持ち、左再帰も扱える多言語対応ツール
\end{itemize}

最後に、パーサーコンビネータという、プログラミング言語の機能を活用して直接パーサーを組み立てる手法を学び、実際にJavaでパーサーコンビネータ「JComb」を実装しました。わずか500行に満たないコードで、PEGに対応した実用的なパーサーコンビネータを作れることを示しました。

構文解析器生成系とパーサーコンビネータは、それぞれ異なる強みを持ちます。構文解析器生成系は宣言的で可読性が高く、パーサーコンビネータはプログラミング言語の型システムやIDEの支援を受けられ、柔軟な拡張が可能です。どちらを選ぶかは、プロジェクトの要件や開発チームの好みによって決めることになるでしょう。
